%% =============================================================
%%   Poster Template
%% =============================================================

\documentclass[final]{beamer}

%% ---------------------- POSTER OPTIONS ----------------------
%% Size
\usepackage[size=a0, scale=1.1]{beamerposter}                                   %% A0 landscape, 118.9 x 84.1 cm (beamerposter's default)
%%\usepackage[size=a0, orientation=portrait, scale=1.1]{beamerposter}           %% A0 portrait (try \jajcolumns{2})
%%\usepackage[size=custom, width=121.92, height=91.44, scale=1.3]{beamerposter} %% 48 x 36 in landscape (common in the US)

%% Theme
\usetheme[institution=uwmadison]{jajposter}

%% Columns
\jajcolumns{3}             %% number of columns
\jajcolumnsep{0.025}       %% gap between columns (fraction of paper width)

%% ---------------------- PACKAGES ----------------------------
\usepackage{jaj-macros}
\usepackage{microtype}
\usepackage{graphicx}
\graphicspath{{figures/}}
\usepackage{booktabs}
\usepackage{siunitx}
\usepackage{algorithm}
\usepackage{algpseudocode}
\usepackage{qrcode}
\usepackage[authoryear,round]{natbib}
\hypersetup{colorlinks=false, pdftitle={Poster Title}, pdfauthor={Jacob A. Johnson}}

%% ---------------------- TITLE INFO --------------------------
\title{Poster Title Goes Here}
\author{Jacob A.\ Johnson\inst{1} \and Coauthor Name\inst{1}}
\institute{\inst{1} Department of Statistics, University of Wisconsin--Madison}

%% ---------------------- LOGOS -------------------------------
%% The institution logo appears on the right by default.
%% \logoright{}                                              %% remove it
%% \logoleft{\includegraphics[height=\jajlogoheight]{lab-logo.png}}
%% \setlength{\jajlogoheight}{6cm}                           %% resize logos

%% ---------------------- FOOTER ------------------------------
\footercontent{%
  \href{mailto:jajohnsonstats@gmail.com}{jajohnsonstats@gmail.com}%
  \hfill
  \href{https://github.com/jajohnsonstats}{github.com/jajohnsonstats}%
}


%% ---------------------- DOCUMENT ----------------------------
\begin{document}
\begin{frame}[t]
\begin{postercolumns}

  %% -------------------- Column 1 ----------------------------
  \begin{column}{\colwidth}

    \begin{block}{Introduction}
      Your introduction text goes here.
      \begin{itemize}
        \item First point, with an \alert{alerted phrase}.
        \item Second point.
      \end{itemize}
    \end{block}

    \begin{block}{Background}
      Background content here.
      \heading{A sub-heading}
      More text.
    \end{block}

  \end{column}
  \separatorcolumn

  %% -------------------- Column 2 ----------------------------
  \begin{column}{\colwidth}

    \begin{block}{Methods}
      Methods go here.
      \begin{theorem}[Optional name]
        A result, typeset inline and numbered.
      \end{theorem}
    \end{block}

    \begin{alertblock}{Key Result}
      Use an alert block for the one thing people should remember.
    \end{alertblock}

    \begin{block}{Results}
      Results go here.
      %% \begin{figure}
      %%   \centering
      %%   \includegraphics[width=0.9\linewidth]{result.pdf}
      %%   \caption{Caption text.}
      %% \end{figure}
    \end{block}

  \end{column}
  \separatorcolumn
  
  %% -------------------- Column 3 ----------------------------
  \begin{column}{\colwidth}

    \begin{block}{Discussion}
      Discussion text here.
    \end{block}

    \begin{block}{Conclusion}
      Conclusion text here.
    \end{block}

    \begin{exampleblock}{Paper \& Code}
      \begin{minipage}[c]{0.68\linewidth}
        Scan for the preprint and code.
      \end{minipage}\hfill
      \begin{minipage}[c]{0.28\linewidth}
        \raggedleft\qrcode[height=0.9\linewidth]{https://github.com/jajohnsonstats}
      \end{minipage}
    \end{exampleblock}

    \begin{exampleblock}{References}
      \footnotesize
      %% \nocite{*}   %% list every entry in references.bib without citing it
      \bibliographystyle{plainnat}
      \bibliography{references}
    \end{exampleblock}

  \end{column}

\end{postercolumns}
\end{frame}
\end{document}
